\section{Testes por camada}
\label{sec:testes-camadas}

A arquitetura acima existe, em boa parte, para tornar os testes da
Parte~\ref{parte:testes} baratos. O capítulo~\ref{cap:ponte-quarkus} apresentou
\texttt{@QuarkusTest}, \texttt{@QuarkusIntegrationTest} e Dev Services;
aqui a pergunta é \emph{onde} usar cada um.

\small
\begin{longtable}{@{}p{3cm} p{4.2cm} p{5cm}@{}}
\toprule
Camada & Tipo de teste & Por quê \\
\midrule \endhead
\texttt{dominio} & JUnit puro, rápido, muitos & Sem framework, sem
  container; o ciclo de TDD da aula~1 roda em milissegundos \\
\texttt{aplicacao} & JUnit puro com fakes das portas & As portas são
  interfaces suas; um \texttt{Pedidos} em memória substitui o banco \\
\texttt{adaptador.\allowbreak saida.*} & \texttt{@QuarkusTest} com Dev Services &
  O que importa é o SQL, o mapeamento e a configuração reais \\
\texttt{adaptador.\allowbreak entrada.\allowbreak rest} & \texttt{@QuarkusTest} pela API HTTP,
  com o caso de uso substituído quando convém & Serialização, códigos de
  status e validação \\
Artefato final & \texttt{@QuarkusIntegrationTest} & Caixa preta contra o
  JAR ou o executável nativo produzido pelo build \\
\bottomrule
\end{longtable}
\normalsize

\paragraph{Substituir um bean num \texttt{@QuarkusTest}.}
O guia \citetitle{doc:quarkus-testing} mostra o \texttt{@InjectMock},
disponível com a dependência \texttt{quarkus-junit-mockito}, e o
\texttt{QuarkusMock}, com \texttt{installMockForType} para instalar um
mock antes dos testes. Na arquitetura hexagonal, o candidato natural a
mock é uma porta (\texttt{CobrancaGateway}), nunca uma classe de
domínio. O \texttt{QuarkusMock} só funciona com beans de escopo normal;
para o \texttt{FecharPedido} \texttt{@Singleton}, use \texttt{@InjectMock}
com \texttt{@MockitoConfig(convertScopes = true)}.

\paragraph{Cuidado com o custo de subir o Quarkus.}
O mesmo guia diz que, nos exemplos dele, o Quarkus sobe uma vez para
todos os testes. Um teste com \texttt{@TestProfile} diferente do anterior
faz o Quarkus ser desligado e iniciado de novo com o novo perfil. Agrupe
os testes por perfil e prefira poucos perfis: cada um soma uma subida ao
tempo da suíte.

\paragraph{Dev Services e o modo nativo.}
Os Dev Services (\citetitle{doc:quarkus-dev-services}) sobem banco e
outros serviços com Testcontainers quando nada foi configurado, nos
modos de desenvolvimento e teste. O guia de testes diz que o
\texttt{@QuarkusIntegrationTest} funciona com contêineres iniciados por
Dev Services, o que permite rodar os mesmos testes de caixa preta contra
o executável nativo no CI. É o teste que pega os problemas da
seção~\ref{sec:reflexao-nativo}: reflexão não registrada só aparece no
binário nativo.

\paragraph{Mutação.}
O PIT (aula~2) brilha em \texttt{dominio} e \texttt{aplicacao}, onde os
testes são rápidos e o código não depende do Quarkus. A integração do
PIT com \texttt{@QuarkusTest} continua não verificada
(capítulo~\ref{cap:ponte-quarkus}); manter as regras fora do framework
evita depender dela.
